\documentclass[11pt,a4paper]{article}
\usepackage[margin=2.6cm]{geometry}
\usepackage{amsmath,amssymb,amsthm}
\usepackage{booktabs}
\usepackage[hidelinks]{hyperref}
\newtheorem{proposition}{Proposition}
\newtheorem{lemma}[proposition]{Lemma}
\title{Exact values and certified lower bounds for the no-three-in-line problem in the cube}
\author{Aleksei Kudriashov (Alex Komang)\\ \small Nusa Dua Studio, Nusa Dua, Bali, Indonesia\\ \small ORCID: 0009-0006-8885-5338}
\date{3 September 2026 --- version 1.5\\ \small\texttt{doi:10.5281/zenodo.22019279}}
\begin{document}\maketitle
\begin{abstract}
Let $a(n)$ be the largest number of points of the grid $\{0,1,\dots,n-1\}^3$ no three of which are collinear.
P\'or and Wood proved $a(n)=\Theta(n^2)$, the lower bound by the construction
$\{(x,y,x^2+y^2\bmod p)\}$, which has no three collinear precisely when $p\equiv3\pmod4$
(Lemma~4 of \cite{PorWood}); no exact values appear to have been recorded, and the sequence was not in
the OEIS when this work began (it is now A399138, approved August 2026).  We determine
\[a(1),\dots,a(6)=1,\,8,\,16,\,28,\,40,\,64,\]
give certified lower bounds $a(7)\ge73$, $a(8)\ge94$, $a(9)\ge116$, $a(10)\ge138$, $a(11)\ge164$ (the bounds for $n\ge8$,
new in version 1.5, come from exact optimisation inside the symmetry classes of the cube group), and prove $a(p)\ge p^2$ for every prime $p$ by an
elementary construction from an anisotropic binary quadratic form.

The method is a certified decision procedure rather than a search: a line carrying at most two lattice points
forbids nothing, so admissibility is exactly ``at most two chosen points on every line with three or more
lattice points'', which is a propositional formula on $n^3$ variables.  Unsatisfiability is answered with a
DRAT certificate that an independent checker verifies, and the lower bounds are witnesses re-verified by
exhaustive integer arithmetic over every triple.  We are explicit about the three separate links this
requires --- that the formula expresses the problem, that it is unsatisfiable, and that the symmetry pruning
discards nothing --- because a certificate attests only to the formula it was given, and an error in the other
two links would pass through it unseen.

Along the way the optima turn out to share a layer structure that fails at exactly one of the values computed:
for $n=2,3,4,6$ the two outer layers carry the planar maximum $2n$ and the inner ones $2n-2$, giving
$2n^2-2n+4$; at $n=5$ two outer layers provably cannot both carry $2n$ at any total, and at $n=7$ the
corresponding profile is impossible as well.  We have no explanation for this.
\end{abstract}

\section{Definitions and elementary bounds}
Points of $G_n:=\{0,\dots,n-1\}^3$; a set is \emph{lawful} if no three of its points are collinear in
$\mathbb R^3$ (equivalently, all $2\times2$ minors of the difference matrix vanish for no triple).
Each of the $3n^2$ axis-parallel lines carries at most two points of a lawful set, and each point lies on
three of them, so $|S|\le 2n^2$; this is the trivial upper bound, attained at $n=2$.

\begin{proposition}\label{p:alg}
For every prime $p$, $a(p)\ge p^2$.
\end{proposition}
\begin{proof}
Let $Q(x,y)=x^2-dy^2$ with $d$ a quadratic non-residue modulo $p$, so that $Q(a,b)\equiv0$ only for
$a\equiv b\equiv0$ (an anisotropic binary form; one exists for every $p$).  Put
$S=\{(x,y,Q(x,y)\bmod p):0\le x,y<p\}$, of size $p^2$.  Suppose three of its points are collinear.  Their
projections to the $(x,y)$-plane are distinct (the third coordinate is a function of the first two), hence lie
on a line with primitive integer direction $(a,b)$, and the three points are $(x_0+t_ia,\,y_0+t_ib,\,\cdot\,)$
with distinct integers $t_i$, $|t_i|<p$.  Along this line the third coordinate satisfies
$z\equiv Q(x_0+ta,y_0+tb)=Q(a,b)\,t^2+\ell(t)\pmod p$ with $\ell$ affine, whereas collinearity in
$\mathbb R^3$ forces $z$ to be an affine function of $t$.  A quadratic congruence agreeing with an affine one
at three points distinct modulo $p$ has vanishing leading coefficient, so $Q(a,b)\equiv0$ and therefore
$(a,b)=(0,0)$, a contradiction.
\end{proof}
Neither $x^2+y^2$ (anisotropic iff $p\equiv3\bmod 4$) nor $x^2+xy+y^2$ (iff $p\equiv2\bmod3$) works for all
$p$; the choice $x^2-dy^2$ does.  Proposition~\ref{p:alg} is thereby the completion, to every odd prime, of
Lemma~4 of P\'or and Wood \cite{PorWood}, whose construction is the case $Q=x^2+y^2$ and works precisely for
$p\equiv3\pmod4$; we verified their dichotomy directly ($6$ and $286$ collinear triples at $p=5$ and $p=13$, none
at $p=7$, $11$).  For $p=2$ the statement is trivial: no line meets $\{0,1\}^3$ in three points, so
$a(2)=8>4=p^2$.  The bound of Proposition~\ref{p:alg} is not sharp: $a(5)\ge40>25$ and
$a(7)\ge73>49$.  As a guard against a slip in transcribing the construction rather than in the proof,
we also built $S$ explicitly for $p=3,5,7,11,13$ and tested \emph{every} one of its triples by integer cross
product --- $790\,244$ triples at $p=13$ --- finding none collinear.

\section{Method: a certified decision procedure}
A line carrying at most two lattice points of $[n]^3$ forbids nothing, so
\begin{equation}\label{eq:reform}
S\subseteq[n]^3 \text{ is admissible}\iff |S\cap L|\le 2 \text{ for every line } L \text{ with } \ge3 \text{ lattice points.}
\end{equation}
This turns the question ``is there an admissible set of size $M$?'' into a Boolean satisfiability problem with
$n^3$ variables --- $125$ for $n=5$, $216$ for $n=6$ --- one at-most-two cardinality constraint per rich line,
and one global at-least-$M$ constraint.  The global counter is a totalizer truncated at $2n^2$; the truncation
is legitimate because each of the $n^2$ axis-parallel lines carries at most two points, so $a(n)\le 2n^2$.
We add the $47$ lexicographic leader constraints of the cube group, which preserve satisfiability because the
property in \eqref{eq:reform} is invariant under that group.

We adopted this route after measuring the alternative.  A prefix-split exhaustive search on $42$ cores closed
$5$ of $204$ prefixes in two hours at $n=5$, extrapolating to roughly $88$ core-days; the same question is
decided by a SAT solver in minutes.  Speed, however, is not the reason.  An exhaustive search establishes
completeness by \emph{its own} covering argument, which only its authors' code can check; an unsatisfiability
answer comes with a DRAT certificate that an independent, widely used checker verifies without trusting our
solver or our code at all.

\medskip\noindent\textbf{Three separate links.}  It is worth being explicit about what is and is not certified,
because the three claims below rest on different foundations.
\begin{enumerate}\itemsep2pt
\item \emph{The formula expresses the problem.}  Verified two independent ways: by comparing the set of
triples forbidden by the rich-line constraints against the set with vanishing cross product, and by forcing each
collinear triple as an assumption and requiring unsatisfiability.  At $n=4$ and $n=5$ there are $376$ and $1858$
collinear triples; none was left unforbidden and none was forbidden spuriously.  A DRAT checker cannot verify
this link --- it knows nothing of cubes and lines --- and this is exactly where an encoding bug would hide.
\item \emph{The formula is unsatisfiable.}  Verified by \texttt{drat-trim} on the certificate emitted by
\texttt{kissat}.
\item \emph{A witness is admissible.}  Verified by exhaustive integer arithmetic over all $\binom{|S|}{3}$
triples, by a program sharing no code with the search.
\end{enumerate}
Links 1 and 2 give the upper bound; links 1 and 3 give the lower bound.

\medskip\noindent\textbf{Removing link 1 for small $n$.}  The two checks above verify that the encoding is
faithful.  A stronger thing is possible: to compute the same values without an encoding at all, so that link 1
does not need to be trusted because it is not used.  We wrote a plain exhaustive traversal that places points
in the cube and tests triples by integer cross product --- the literal definition --- with no propositional
formula, no solver and no certificate.  Two points kill the line through them, so every cell carries a counter
of dead lines covering it and a cell is available while that counter is zero; a branch is abandoned when the
points placed plus the cells still available cannot exceed the best found.  Run from a bar of zero, so that no
value can be fed in and echoed back:
\begin{center}
\begin{tabular}{rrr}
\toprule
$n$ & result & nodes \\
\midrule
1 & $1$ (exhausted) & $2$\\
2 & $8$ (exhausted) & $37$\\
3 & $16$ (exhausted) & $49\,812$\\
\bottomrule
\end{tabular}
\end{center}
These agree with the certified values and share nothing with the certified route but the definition of the
prohibition itself.  Our own full enumeration at $n=3$ visits all $3\,813\,884$ lawful subsets (the empty one
included) and finds exactly \emph{two} optima of size $16$ --- a chiral pair, each with stabiliser of order $24$
in the full cube group, both with layer profile $[6,4,6]$ along all three axes.  At $n=4$ a traversal from a bar
of $27$ \emph{exhausted} its tree ($9\,012\,859\,649$ nodes, $3180$ s), proving $a(4)=28$ with no encoding at
all --- a bar is harmless under exhaustion, since it prunes only branches unable to exceed it.  The count at $n=4$
is now ours as well: a $32$-shard traversal (sharded by first cell, run to completion on every shard,
$13\,215\,017\,061$ nodes in all) finds exactly $14$ optima of size $28$, in three classes under the full cube
group with stabiliser orders $8$, $8$, $24$ ($6+6+2=14$) --- in agreement with an independent external recount.
Every optimum at $n\le4$ therefore has the layer profile $[2n,(2n-2)^{n-2},2n]$: the empirical observation of
Section~5 is a fact up to $n=4$.

\emph{Correction.}  An earlier version of this paragraph quoted two numbers received from the first solver's
message and never passed through the verification bench of this very note: ``$2\,220\,075$ subsets'' and ``the
count $10\,960$''.  Both are withdrawn, and both origins have since been reconstructed: $2\,220\,075$ is exactly
$\binom{27}{8}$ --- the enumeration of all $8$-subsets of $3^3$ that verifies the uniqueness of the $8$-point
optimum of the \emph{no-four-coplanar} problem (A280538 gives $1$ there) --- and $10\,960=48\cdot232-176$ is
the number of labelled optima of that same neighbouring problem at $n=4$, where A280538 gives $232$ classes.
Two numbers from the wrong problem, quoted into the section about removing trust; recorded, not smoothed.

We are explicit about how far this extends.  From a bar of zero the traversal reaches $28$ at $n=4$ without
exhausting in an hour ($8.85\cdot10^9$ nodes); from a bar of $27$ it exhausts ($9.01\cdot10^9$ nodes, $53$
minutes), so $a(4)=28$ is established twice over as well.  The honest division: for $n\le4$ both the value and
the bound are encoding-free in addition to certified; for $n\ge5$ the encoding is checked but used.

\medskip\noindent\textbf{Calibration.}  Before computing anything new we recomputed, as a calibration, the two
smallest nontrivial values $a(3)=16$ and $a(4)=28$ --- there is no prior published source for them that we could
find, so the calibration is against our own earlier runs, not against the literature --- each with unsatisfiability
at $M+1$ in under a second, independently on two machines.
For instances that resist a single solver call we split into independent pieces by fixing the contents of the
first columns; the split is complete because three points of a column are collinear, so subsets of size at most
two exhaust it.  The aggregator refuses a verdict unless every piece of the manifest is present and
unsatisfiable --- a missing piece and an unsatisfiable one must never look alike.

\section{Exact values}
\begin{center}\begin{tabular}{rrrl}\toprule
$n$ & $a(n)$ & $a(n)/n^2$ & how established\\\midrule
1 & 1 & --- & trivial\\
2 & 8 & 2.00 & exhaustive; meets the trivial bound $2n^2$\\
3 & 16 & 1.78 & exhaustive, and certified\\
4 & 28 & 1.75 & exhaustive, and certified\\
5 & 40 & 1.60 & certified\\
6 & 64 & 1.78 & certified\\\bottomrule
\end{tabular}\end{center}
\noindent The value $a(6)=64$ appears to be new; a search of the OEIS for $1,8,16,28,40$ and for
$8,16,28,40,64$ returns nothing, and the sequence itself is not in the database.

For $n\le4$ the values were also produced by two independent exhaustive programs whose node counts differ by
a factor of about $72$ ($17/3855/4.21\cdot10^7$ against $17/20068/3.03\cdot10^9$ for $n=2,3,4$), so their
agreement is not an artefact of a shared design.

For $n=5$ and $n=6$ the decisive question --- does a set of $41$, respectively $65$, points exist? --- was
split into independent pieces by fixing the contents of the first two $z$-columns, and the heaviest pieces
were split once more on the third column.  Both splits are complete because three points of a column are
collinear, so subsets of size at most two exhaust it.  The final tally, reproduced in
\texttt{logs/no3\_3d/FINAL\_VERDICT.txt}:
\begin{center}\begin{tabular}{lrrr}\toprule
 & $n=5$, $M=41$ & $n=6$, $M=65$\\\midrule
top-level pieces & $256$ & $484$\\
closed directly & $221$ & $450$\\
closed through all their continuations & $35$ & $34$\\
continuations required per piece & $16$ & $22$\\
satisfiable pieces found & $0$ & $0$\\
pieces left undecided & $0$ & $0$\\\bottomrule
\end{tabular}\end{center}
The aggregator builds the expected set of pieces from the manifest rather than from the results, treats a
missing piece and an undecided one as failures, and distinguishes three outcomes rather than two: only an
explicit unsatisfiable answer closes a piece, a satisfiable one anywhere would destroy the claim, and a
killed or timed-out run counts as \emph{no information} --- it neither closes nor refutes.  Eleven such
uninformative lines occur in the $n=6$ record, from runs interrupted when work was moved between machines;
each of the pieces concerned is closed by an explicit unsatisfiable answer from another run.

The lower bounds are witnesses, checked by exhaustive integer arithmetic over all $\binom{|S|}{3}$ triples by
a program sharing no code with the search: $40$ points at $n=5$ ($9880$ triples) and $64$ points at $n=6$
($41664$ triples).  For $n=6$ two witnesses were obtained independently, by a SAT solver and by a constraint
solver in decision form; they turned out to be the same set, which the stabiliser explains --- it has order
$24$, so the configuration has only two distinct images under the $48$ symmetries of the cube.

As a further check that is not part of the proof, a constraint solver was given the $n=6$, $M=65$ question in
one piece for an hour and returned no answer.  That is not evidence of impossibility --- it is the solver
saying it does not know --- but it found no counterexample either, while at $M=64$ it produces a solution in
seconds.

\section{Structure of the optima, and an anomaly at $n=5$}
Writing the layer profile of a set as the number of its points in each of the $n$ layers perpendicular to an
axis, every optimum we have found has the \emph{same} profile along all three axes:
\begin{center}\begin{tabular}{rrl}\toprule
$n$ & $a(n)$ & layer profile (identical for all three axes)\\\midrule
2 & 8  & $4,4$\\
3 & 16 & $6,4,6$\\
4 & 28 & $8,6,6,8$\\
5 & 40 & $8,8,8,8,8$\\
6 & 64 & $12,10,10,10,10,12$\\\bottomrule
\end{tabular}\end{center}
A layer is an $n\times n$ grid with no three collinear, so it holds at most $2n$ points.  For $n=2,3,4,6$ the two
outer layers attain that planar maximum and the inner ones hold $2n-2$, giving $2\cdot 2n+(n-2)(2n-2)=2n^2-2n+4$
--- which reproduces $8,16,28$ and $64$ exactly.

At $n=5$ the pattern fails, and it fails for a structural reason rather than by accident.  Prescribing the
profile as a hard constraint and deciding feasibility shows that
\begin{center}\begin{tabular}{ll}\toprule
profile at $n=5$ & verdict\\\midrule
$10,8,8,8,10$ (total $44$, the value the pattern predicts) & infeasible\\
$10,9,8,9,10$; $10,8,9,8,10$; $10,10,10,10,10$; $9,9,9,9,9$ & infeasible\\
$10,8,8,8,8$ (total $42$) and $10,7,8,7,10$ (total $42$) & infeasible\\
$\mathbf{10,6,8,6,10}$ \textbf{(total }$\mathbf{40}$\textbf{)} & \textbf{infeasible}\\
$8,8,8,8,8$ (total $40$) & feasible\\\midrule
$n=6$: $12,10,10,10,10,12$ (total $64$) & feasible\\
$n=7$: $14,12,12,12,12,12,14$ (total $88$) & infeasible\\\bottomrule
\end{tabular}\end{center}
A caveat on how these verdicts are to be read.  Prescribing the profile along all three axes at once is a
restriction, so an infeasibility obtained that way would exclude only symmetric configurations, not that number
of points.  Every infeasibility above was therefore re-derived with the profile prescribed along a
\emph{single} axis, which excludes every configuration with that layer distribution.  The satisfiable rows need
no such care: a solution is a solution.
The fourth line is the sharpest: at $n=5$ two outer layers cannot both attain the planar maximum $2n$ at
\emph{any} total, so the optimum is forced onto a uniform profile.  The same computation at $n=6$ finds the
analogous profile realisable, and at $n=7$ it is again impossible.  The pattern $[2n,(2n-2)^{n-2},2n]$ is thus
realisable at $n=2,3,4,6$ and impossible at $n=5,7$ --- the obstruction is not an artefact of one case, and it
is not an obstruction to the idea of maximal outer layers as such.

The obstruction can be priced exactly.  A layer is a planar no-three-in-line configuration and so carries at
most $2n$ points; requiring \emph{both} outer layers to attain that planar maximum, and maximising the total,
gives
\begin{center}\begin{tabular}{rrrl}\toprule
$n$ & both outer layers & maximum total & $a(n)$\\\midrule
4 & $8$ & $28$ & $28$ --- costs nothing\\
5 & $10$ & $\mathbf{39}$ & $\mathbf{40}$ --- \textbf{costs exactly one point}\\
6 & $12$ & $64$ & $64$ --- costs nothing\\\bottomrule
\end{tabular}\end{center}
So the anomaly is not that a fitted formula occasionally fails.  It is that two planar-optimal configurations
placed in the outer planes are compatible with the spatial optimum at $n=4$ and $n=6$ and incompatible at
$n=5$, and incompatible by a single point.

We do not have a proof of why.  A count of the planar configurations attaining $2n$ points suggests rigidity:
the $4\times4$ grid admits $11$ such configurations in $4$ classes under the square's symmetries, the $5\times5$
grid admits $32$ in only $5$ classes, and the $6\times6$ grid admits $50$ in $11$ classes.  We offer this as an
observation accompanied by its data, not as an argument.  One natural mechanism can be ruled out: the midpoint
of a segment joining the two outer layers is a lattice point when $n$ is odd, and must stay empty, but an exact
count shows the middle layer still retains up to $16$ free cells --- above its own planar ceiling of $10$ --- so
that constraint does not bind.

\section{Certified lower bounds for $n=7,\dots,11$}
\begin{center}\begin{tabular}{rrrl}\toprule
$n$ & $a(n)\ge$ & $/n^2$ & source of the witness\\\midrule
7 & 73 & 1.49 & fourteen inequivalent configurations known; the first was invariant under every permutation of the axes (order 6)\\
8 & 94 & 1.47 & two inequivalent configurations, each invariant under a subgroup of order 4 (version 1.4 had $93$, order 6)\\
9 & 116 & 1.43 & two inequivalent configurations: stabilisers of order 4 and of order 12; the latter is optimal within its class\\
10 & 138 & 1.38 & full stabiliser of order 12; optimal within the order-12 class\\
11 & 164 & 1.36 & stabiliser of order 12; the bound inside the class is not closed\\\bottomrule
\end{tabular}\end{center}
The bounds for $n=8,\dots,11$ (version 1.5, 2--3 September 2026) come from a sweep over all $33$ conjugacy
classes of subgroups $H$ of the symmetry group of the cube (order $48$): in each class the maximum of $|S|$
over $H$-invariant admissible sets is computed by exact optimisation (CP-SAT over the $H$-orbits of cells, with
``at most two chosen cells on every grid line'') under a time limit of $25$--$40$ minutes per class.  Where the
solver closes the bound the value is the exact maximum of the class (this happened for the order-$12$ class at
$n=9$ and $n=10$); elsewhere it is a lower bound inside the class.  Every configuration was verified by two
independent programs over all $\binom{|S|}{3}$ triples, its stabiliser computed under the $48$ symmetries, and
its local optimality certified: no point can be replaced (all are rigid, min $\kappa\ge2$ in the sense of the
disjoint-killers lemma of the companion repository \texttt{lemma-atelier}), and no exchange removing at most
four points ($73$, $94$, $116$) or three points ($138$, $164$) improves it.  The configurations, the per-class
results and the verifier are archived separately \cite{Witnesses}.

The rest of this section is the version-1.4 account of $n=7,8$, kept as data.  Each configuration is re-verified by testing all $\binom{|S|}{3}$ triples with integer cross products.  These
are lower bounds only; we make no claim about how far they are from the truth, and we expect them to be weak,
since a randomised search is a poor instrument at this size.  The $n=7$ bound illustrates the point: it stood
at $71$ from randomised search and rose to $73$ within the hour once the search was restricted to
configurations invariant under a cyclic permutation of the axes.  Restricting to a subgroup collapses the cells
into orbits and shrinks the search by the order of the group; it is legitimate for lower bounds only, since a
configuration found is a configuration, while failure to find one among the symmetric ones says nothing about
the rest.  The size of the group has an optimum in the middle rather than at either end, and the optimum is a balance:
the group must be large enough to make the problem solvable and small enough that good configurations still
fall inside the symmetric stratum.  At $n=8$:
\begin{center}\begin{tabular}{rl}\toprule
order of the group & best found\\\midrule
2 & 82, 84\\ 3 & 88\\ 4 & 88, 90\\ \textbf{6} & \textbf{92, 92, 93}\\ 8 & 80, 88\\ 12 & 88\\
24 & 88 (proved optimal within its stratum)\\ 48 & 80 (proved optimal within its stratum)\\\bottomrule
\end{tabular}\end{center}
(The full sweep of version 1.5, $40$ minutes per class, later gave $94$ in two classes of order $4$, so the
order-$6$ value $93$ of this table was not the maximum of its own neighbourhood of classes either.)
At $|G|=48$ only $20$ orbits remain, the solver settles the stratum in seconds, and the answer is $80$, because
genuinely good configurations do not have that much symmetry.  Note also that the order alone does not decide
it: two different groups of order $8$ gave $80$ and $88$.  We record this as data.  An earlier draft of this
note stated instead that a larger group simply gives a better result --- a pattern read off the first four
rows, which the remaining four contradict.  (An earlier draft of this note quoted
$72$ and $90$ here.  Those numbers came from a run whose witness was never saved, and when we came to attach
witnesses to them the journals we could reproduce gave $71$ and $89$.  We report the values we can exhibit.)
Two indications.  First, the layer structure of
the smaller optima suggests $2n^2-2n+4=88$ at $n=7$ --- and although we have shown that the corresponding
\emph{profile} is impossible, that leaves the value itself untouched.  Second, at $n=6$ the same randomised
search reached $64$, which the certified computation then confirmed to be optimal, so the instrument is not
hopeless either; we simply do not know which case $n=7$ resembles.

For the record we also measured what the abandoned route would have cost, though the measurement deserves a
caveat.  Knuth's random-probe estimator, calibrated against the exact tree at $n=4$, put an unaided exhaustive
proof at $n=6$ at some $4.6\cdot10^{22}$ nodes, and that figure is what sent us to a decision procedure.  We
later found the plain estimator unreliable on trees of this family --- on a sibling problem the mean of twenty
thousand probes was $89\%$ a single probe, with the median five orders of magnitude below the mean --- and
replaced it by a stratified version that agrees with the exact tree at $n=4$ to within $0.1\%$.  The
$4.6\cdot10^{22}$ was produced by the plain estimator and has not been recomputed with the stratified one, so
it should be read as an order of magnitude and nothing finer.  It says nothing about the difficulty of the
problem in any case, only about the difficulty of one way of attacking it.

\section{Reproducibility, and what is \emph{not} claimed}
All programs, journals, witnesses, manifests and the verification protocol are public:
\url{https://github.com/iwasborninbali/saturation}.  The protocol \texttt{docs/VERIFICATION.md} is written for
a reader who trusts none of our code and gives the commands in order.

The following are the limits of what is established here, stated so that a reader need not infer them.
\begin{itemize}\itemsep2pt
\item At $n=5$ the upper bound no longer rests on one solver: the same $256$-piece split was closed
independently by Glucose, of the MiniSat lineage rather than kissat's, with every piece unsatisfiable and none
satisfiable, the last taking $2785$ seconds and $14.4$ million conflicts.  Two solvers of unrelated descent, two
independently written aggregators, two independent semantic checks and two independently verified witnesses now
stand behind $a(5)=40$.  Worth recording separately: the \emph{monolithic} instance defeated Glucose for
$10{,}000$ seconds while its pieces fell in minutes, so the case split is a property of the method rather than
of one solver.  \textbf{At $n=6$ this has not been done} --- there are $484$ pieces and they are heavier --- and
the reader should not carry $n=5$ over to it.
\item The upper bound for $n=6$ rests on one SAT solver.  The instances that a single call decides carry
DRAT certificates verified by \texttt{drat-trim}, which removes the dependence on the solver; the instances
that had to be split do not carry stored certificates, because the certificates run to tens of gigabytes.
They are regenerated deterministically on demand, and the case split, its completeness and its aggregation are
each verified separately.
\item We have no proof of the upper bounds independent of the propositional encoding.  Our own exhaustive
enumerator does not reach $n=5$: it closed $24$ of $4096$ prefixes at $15$--$33$ billion nodes each, with the
full traversal on the order of $10^{14}$ nodes.  It corroborates the lower bound only.
\item The verification of the symmetry pruning shows that our \emph{implementation} discards nothing; it is
not a proof of the underlying lex-leader argument, which is standard but which we did not reprove.
\item The stratum sweep (version 1.5) proves nothing about $a(n)$: a class maximum is an upper bound for that
class only, and the classes are enriched in maxima but need not contain them --- at $n=6$ every class with a
non-trivial symmetry stays below $64$ except those attained by the certified optimum itself.
\item The layer-structure anomaly at $n=5$ and $n=7$ is a computation, not an explanation.  A count of the
planar configurations attaining $2n$ points suggests rigidity ($11$ configurations in $4$ classes for
$4\times4$; $32$ in only $5$ classes for $5\times5$; $50$ in $11$ classes for $6\times6$), and one natural
mechanism --- the lattice midpoint between two outer layers --- is ruled out by an exact count.  Beyond that we
do not know.
\end{itemize}

\begin{thebibliography}{9}
\bibitem{PorWood} A.~P\'or and D.~R.~Wood, No-Three-in-Line-in-3D, \emph{Algorithmica} 47 (2007) 481--488;
preliminary version in Proc.\ Graph Drawing 2004, Lecture Notes in Comput.\ Sci.\ 3383, Springer.
\bibitem{Kissat} A.~Biere, K.~Fazekas, M.~Fleury and M.~Heisinger, CaDiCaL, Kissat, Paracooba, Plingeling and
Treengeling entering the SAT Competition 2020, Proc.\ SAT Competition 2020, University of Helsinki, 2020.
\bibitem{DratTrim} N.~Wetzler, M.~J.~H.~Heule and W.~A.~Hunt, Jr., DRAT-trim: efficient checking and trimming
using expressive clausal proofs, Proc.\ SAT 2014, Lecture Notes in Comput.\ Sci.\ 8561, Springer, 2014.
\bibitem{Glucose} G.~Audemard and L.~Simon, Predicting learnt clauses quality in modern SAT solvers,
Proc.\ IJCAI 2009.
\bibitem{OEIS} OEIS Foundation Inc., The On-Line Encyclopedia of Integer Sequences, \url{https://oeis.org},
sequence A399138.
\bibitem{Witnesses} A.~Kudriashov, Witness configurations and lower bounds for the no-three-in-line problem in the
$n\times n\times n$ grid (OEIS A399138), Zenodo, 2026, \url{https://doi.org/10.5281/zenodo.22271375}; source
\url{https://github.com/iwasborninbali/a399138-witnesses}.
\end{thebibliography}

\section*{AI/tool-use disclosure}
The searches, the programs and the text were produced by autonomous AI agents (Anthropic Claude) working under
the direction of the author of record, who posed the question, chose what to compute, decided what to claim and
takes responsibility for the content.  Two agents worked with deliberately different program designs and audited
each other; several errors reported in the repository were caught that way and not by their author, including
one --- an unverified symmetry pruning --- that no certificate, witness check or coverage check would have
detected.

The work is arranged so that this should not matter.  Every claim above is meant to be checkable by a third
party without trusting the agents or the author: certificates by \texttt{drat-trim}, witnesses by brute-force
integer arithmetic, coverage by a mechanical aggregator that refuses a verdict on a missing or undecided piece,
and the encoding by comparison against the direct collinearity criterion.  Where a claim is not checkable in
that way, it is listed in the preceding section.

\end{document}
